\documentclass[UTF8,11pt]{ctexart}
\usepackage[a4paper,margin=22mm]{geometry}
\usepackage{graphicx,enumitem,amssymb,tabularx}
\usepackage[normalem]{ulem}
\usepackage{xeCJKfntef}
\setlength{\parindent}{0pt}
\setlength{\parskip}{.6em}
\sloppy
\begin{document}
\section*{2018年6月英语四级真题(第1套)}
\subsection*{\textbf{2018} 年\textbf{6} 月大学英语四级真题（第\textbf{1} 套\textbf{)}}


\subsection*{\textbf{Part I} \textbf{Writing} （\textbf{30 minutes}）}


\textbf{Directions:}\textit{ For this part, you are allowed 30minutes to write a short essay on}\textbf{\textit{ the importance}}


\textbf{\textit{of reading ability and how to develop it.}}\textit{ You should write at least 120 words but no more than} \textit{180 words.} \par\noindent\rule{.95\linewidth}{.4pt}\par  \par\noindent\rule{.95\linewidth}{.4pt}\par  \par\noindent\rule{.95\linewidth}{.4pt}\par


\subsection*{\textbf{Part II} \textbf{Listening Comprehension} \textbf{(25 minutes)}}


\subsection*{\textbf{Section A}}


\textbf{Directions:}\textit{ In this section, you will hear three news reports. At the end of each news report, you}


\textit{will hear two or three questions. Both the news report and the questions will be spoken only} \textit{once. After you hear a question, you must choose the best answer from the four choices} \textit{marked A), B), C) and D). Then mark the corresponding letter on}\textbf{\textit{ Answer Sheet 1}}\textit{ with a} \textit{single line through the centre.}


\textbf{Questions 1 and 2 are based on the news report you have just heard.}


1.


\begin{itemize}[leftmargin=2em,itemsep=.2em,topsep=.4em]\item[A.] The return of a bottled message to its owner's daughter.
\item[B.] A New Hampshire man's joke with friends on his wife.
\item[C.] A father's message for his daughter.
\item[D.] The history of a century-old motel.
\end{itemize}

2.


\begin{itemize}[leftmargin=2em,itemsep=.2em,topsep=.4em]\item[A.] She wanted to show gratitude for his kindness.
\item[B.] She wanted to honor her father's promise.
\item[C.] She had been asked by her father to do so.
\item[D.] She was excited to see her father's handwriting.
\end{itemize}

\textbf{Questions 3 and 4 are based on the news report you have just heard.}


3.


\begin{itemize}[leftmargin=2em,itemsep=.2em,topsep=.4em]\item[A.] People were concerned about the number of bees.
\item[B.] Several cases of Zika disease had been identified.
\item[C.] Two million bees were infected with disease.
\item[D.] Zika virus had destroyed some bee farms.
\end{itemize}

4.


\begin{itemize}[leftmargin=2em,itemsep=.2em,topsep=.4em]\item[A.] It apologized to its customers.
\item[B.] It was forced to kill its bees.
\item[C.] It lost a huge stock of bees.
\item[D.] It lost 2.5 million dollars.
\end{itemize}

\textbf{Questions 5 to 7 are based on the news report you have just heard.}


5.


\begin{itemize}[leftmargin=2em,itemsep=.2em,topsep=.4em]\item[A.] It stayed in the air for about two hours.
\item[B.] It took off and landed on a football field.
\item[C.] It proved to be of high commercial value.
\item[D.] It made a series of sharp turns in the sky.
\end{itemize}

6.


\begin{itemize}[leftmargin=2em,itemsep=.2em,topsep=.4em]\item[A.] Engineering problems.
\item[B.] The air pollution it produced.
\end{itemize}

1


\begin{itemize}[leftmargin=2em,itemsep=.2em,topsep=.4em]\item[C.] Inadequate funding.
\item[D.] The opposition from the military.
\end{itemize}

7.


\begin{itemize}[leftmargin=2em,itemsep=.2em,topsep=.4em]\item[A.] It uses the latest aviation technology.
\item[B.] It flies faster than a commercial jet.
\item[C.] It is a safer means of transportation.
\item[D.] It is more environmentally friendly.
\end{itemize}

\subsection*{\textbf{Section B}}


\textbf{Directions:}\textit{ In this section, you will hear two long conversations. At the end of each}


\textit{conversation, you will hear four questions. Both the conversation and the questions will bespo} \textit{ken only once. After you hear a question, you must choose the best answer from the four} \textit{choices marked A), B), C) and D). Then mark the corresponding letter}\textbf{\textit{ on Answer Sheet 1}}\textit{with} \textit{a single line through the centre.}


Questions 8 to 11 are based on the conversation you have just heard.


8.


\begin{itemize}[leftmargin=2em,itemsep=.2em,topsep=.4em]\item[A.] It seems a depressing topic.
\item[B.] It sounds quite alarming.
\item[C.] It has little impact on our daily life.
\item[D.] It is getting more serious these days.
\end{itemize}

9.


\begin{itemize}[leftmargin=2em,itemsep=.2em,topsep=.4em]\item[A.] The man doesn't understand Spanish.
\item[B.] The woman doesn't really like dancing.
\item[C.] They don't want something too noisy.
\item[D.] They can't make it to the theatre in time.
\end{itemize}

10.


\begin{itemize}[leftmargin=2em,itemsep=.2em,topsep=.4em]\item[A.] It would be more fun without Mr. Whitehead hosting.
\item[B.] It has too many acts to hold the audience's attention.
\item[C.] It is the most amusing show he has ever watched.
\item[D.] It is a show inappropriate for a night of charity.
\end{itemize}

\noindent\begin{tabularx}{\linewidth}{@{}lXXXX@{}}11. & \textbf{A.} Watch a comedy. & \textbf{B.} Go and see the dance. & \textbf{C.} Book the tickets online. & \textbf{D.} See a film with the man.\\\end{tabularx}\par

\textbf{Questions 12 to 15 are based on the conversation you have just heard.}


12.


\begin{itemize}[leftmargin=2em,itemsep=.2em,topsep=.4em]\item[A.] Most of her schoolmates are younger than she is.
\item[B.] She simply has no idea what school to transfer to.
\item[C.] There are too many activities for her to cope with.
\item[D.] She worries she won't fit in as a transfer student.
\end{itemize}

13.


\begin{itemize}[leftmargin=2em,itemsep=.2em,topsep=.4em]\item[A.] Seek advice from senior students.
\item[B.] Pick up some meaningful hobbies.
\item[C.] Participate in after-school activities.
\item[D.] Look into what the school offers.
\end{itemize}

14.


\begin{itemize}[leftmargin=2em,itemsep=.2em,topsep=.4em]\item[A.] Give her help whenever she needs it.
\item[B.] Accept her as a transfer student.
\item[C.] Find her accommodation on campus.
\item[D.] Introduce her to her roommates.
\end{itemize}

15.


\begin{itemize}[leftmargin=2em,itemsep=.2em,topsep=.4em]\item[A.] She has interests similar to Mr. Lee's.
\item[B.] She has become friends with Catherine.
\end{itemize}

2


\begin{itemize}[leftmargin=2em,itemsep=.2em,topsep=.4em]\item[C.] She has chosen the major Catherine has.
\item[D.] She has just transferred to the college.
\end{itemize}

\subsection*{\textbf{Section C}}


\textbf{Directions:}\textit{ In this section, you will hear three passages. At the end of each passage, you will}


\textit{hear three or four questions. Both the passage and the questions will be spoken only once.Afte}


\textit{r you hear a question, you must choose the best answer from the four choices marked A),B), C)}


\textit{and D). Then mark the corresponding letter}\textbf{\textit{ on Answer Sheet 1}}\textit{ with a single line}


\textit{through the centre}.


\textbf{Questions 16 to 18 are based on the passage you have just heard.}


16.


\begin{itemize}[leftmargin=2em,itemsep=.2em,topsep=.4em]\item[A.] To investigate how being overweight impacts on health.
\item[B.] To find out which physical drive is the most powerful.
\item[C.] To discover what most mice like to eat.
\item[D.] To determine what feelings mice have.
\end{itemize}

\noindent\begin{tabularx}{\linewidth}{@{}lXXXX@{}}17. & \textbf{A.} When they are hungry. & \textbf{B.} When they are thirsty. & \textbf{C.} When they smell food. & \textbf{D.} When they want company.\\\end{tabularx}\par

18.


\begin{itemize}[leftmargin=2em,itemsep=.2em,topsep=.4em]\item[A.] They search for food in groups.
\item[B.] They are overweight when food is plenty.
\item[C.] They prefer to be with other mice.
\item[D.] They enjoy the company of other animals.
\end{itemize}

\textbf{Questions 19 to 21 are based on the passage you have just heard.}


19.


\begin{itemize}[leftmargin=2em,itemsep=.2em,topsep=.4em]\item[A.] Its construction started before World War
\item[B.] Its construction cost more than \$ 40 billion.
\item[C.] It is efficiently used for transport.
\item[D.] It is one of the best in the world.
\item[I.] 
\end{itemize}

20.


\begin{itemize}[leftmargin=2em,itemsep=.2em,topsep=.4em]\item[A.] To improve transportation in the countryside.
\item[B.] To move troops quickly from place to place.
\item[C.] To enable people to travel at a higher speed.
\item[D.] To speed up the transportation of goods.
\end{itemize}

\noindent\begin{tabularx}{\linewidth}{@{}lXXXX@{}}21. & \textbf{A.} In the 1970s. & \textbf{B.} In the 1960s. & \textbf{C.} In the 1950s. & \textbf{D.} In the 1940s.\\\end{tabularx}\par

\textbf{Questions 22 to 25 are based on the passage you have just heard.}


\noindent\begin{tabularx}{\linewidth}{@{}lXXXX@{}}22. & \textbf{A.} Chatting while driving. & \textbf{B.} Messaging while driving. & \textbf{C.} Driving under age. & \textbf{D.} Speeding on highways.\\\end{tabularx}\par

23.


\begin{itemize}[leftmargin=2em,itemsep=.2em,topsep=.4em]\item[A.] A gadget to hold a phone on the steering wheel.
\item[B.] A gadget to charge the phone in a car.
\item[C.] A device to control the speed of a vehicle.
\item[D.] A device to ensure people drive with both hands.
\end{itemize}

24.


\begin{itemize}[leftmargin=2em,itemsep=.2em,topsep=.4em]\item[A.] The car keeps flashing its headlights.
\end{itemize}

3


\begin{itemize}[leftmargin=2em,itemsep=.2em,topsep=.4em]\item[B.] The car slows down gradually to a halt.
\item[C.] They are alerted with a light and a sound.
\item[D.] They get a warning on their smart phone.
\end{itemize}

\noindent\begin{tabularx}{\linewidth}{@{}lXXXX@{}}25. & \textbf{A.} Installing a camera. & \textbf{B.} Using a connected app. & \textbf{C.} Checking their emails. & \textbf{D.} Keeping a daily record\\\end{tabularx}\par

\subsection*{\textbf{Part} III \textbf{Reading Comprehension} \textbf{(40 minutes)}}


\subsection*{\textbf{Section A}}


\textbf{Directions:}\textit{ In this section, there is a passage with10 blanks. You are required to select one word}


\textit{for each blank from a list of choices given in a word bank following the passage. Read the}


\subsection*{\textit{passage through carefully before making your choices. Each choice in the bank is identified}}


\textit{by a letter. Please mark the corresponding letter for each item on Answer Sheet 2 with a} \textit{single line through the centre. You may not use any of the words in the bank more than once.}


An office tower on Miller Street in Manchester is completely covered in solar panels. They are used to create some of the energy used by the insurance company inside. When the tower was first \_\_26\_\_ in 1962, it was covered with thin square stones. These small square stones became a problem for the building and continued to fall off the face for 40 years until a major renovation was \_\_27\_\_ . During this renovation the building's owners, CIS, \_\_28\_\_ the solar panel company, Solar century. They agreed to cover the entire building in solar panels. In 2004, the completed CIS tower became Europe's largest \_\_29\_\_ of vertical solar panels. A vertical solar project on such a large \_\_30\_\_ has never been repeated since.


Covering a skyscraper with solar panels had never been done before, and the CIS tower was chosen as one of the "10 best green energy projects". For a long time after this renovation project, it was the tallest building in the United Kingdom, but it was \_\_31\_\_ overtaken by the Mill bank Tower.


Green buildings like this aren't \_\_32\_\_ cost-efficient for the investor, but it does produce much less pollution than that caused by energy \_\_33\_\_ through fossil fuels. As solar panels get \_\_34\_\_ , the world is likely to see more skyscrapers covered in solar panels, collecting energy much like trees do. Imagine a world where building the tallest skyscraper wasn't a race of \_\_35\_\_ , but rather one to collect the most solar energy.


\begin{itemize}[leftmargin=2em,itemsep=.2em,topsep=.4em]\item[A.] cheaper
\item[B.] cleaner
\item[C.] collection
\item[D.] competed
\item[E.] constructed
\item[F.] consulted
\item[G.] dimension
\item[H.] discovered
\item[I.] eventually
\item[J.] height
\item[K.] necessarily
\item[L.] production
\item[M.] range
\item[N.] scale
\item[O.] undertaken
\end{itemize}

\subsection*{\textbf{Section B}}


\textbf{Directions:}\textit{ In this section, you are going to read a passage with ten statements attached to} \textit{it. Each statement contains information given in one of the paragraphs. Identify the} \textit{paragraph from which the information is derived. You may choose a paragraph more than} \textit{once. Each paragraph is marked with a letter. Answer the questions by marking the} \textit{corresponding letter on Answer Sheet 2.}


4


\textbf{Some College Students Are Angry That They Have to Pay to Do Their Homework}


A) Digital learning systems now charge students for access codes needed to complete coursework, take quizzes, and turn in homework. As universities go digital, students are complaining of a new hit to their finances that's replacing—and sometimes joining—expensive textbooks: pricey online access codes that are required to complete coursework and submit assignments.


B) The codes—which typically range in price from \$ 80 to \$ 155 per course—give students online access to systems developed by education companies like McGraw Hill and Pearson. These companies, which long reaped big profits as textbook publishers, have boasted that their new online offerings, when pushed to students through universities they partner with, represent the future of the industry.


C) But critics say the digital access codes represent the same profit-seeking ethos (观念) of the textbook business, and are even harder for students to opt out of. While they could once buy second-hand textbooks, or share copies with friends, the digital systems are essentially impossible to avoid.


D) "When we talk about the access code we see it as the new face of the textbook monopoly (垄断), a new way to lock students around this system," said Ethan Senack, the higher education advocate for the U.S. Public Interest Research Group, to BuzzFeed News. "Rather than \$250 (for a print textbook) you're paying \$ 120," said Senack. "But because it's all digital it eliminates the used book market and eliminates any sharing and because homework and tests are through an access code, it eliminates any ability to opt out."


E) Sarina Harpet, a 19-year-old student at Virginia Tech, was faced with a tough dilemma when she first started college in 2015—pay rent or pay to turn in her chemistry homework. She told BuzzFeed News that her freshman chemistry class required her to use Connect, a system provided by McGraw Hill where students can submit homework, take exams and track their grades. But the code to access the program cost \$ 120—a big sum for Harper, who had already put down \$ 450 for textbooks, and had rent day approaching.


F) She decided to wait for her next work-study paycheck, which was typically \$ 150- \$ 200, to pay for the code. She knew that her chemistry grade may take a dive as a result. "It's a balancing act," she said. "Can I really afford these access codes now?" She didn't hand in her first two assignments for chemistry, which started her out in the class with a failing grade.


G) The access codes may be another financial headache for students, but for textbook businesses, they're the future. McGraw Hill, which controls 21\% of the higher education market, reported in March that its digital content sales exceeded print sales for the first time in 2015. The company said that 45\% of its \$ 140 million revenue in 2015 "was derived from digital products."


H) A Pearson spokesperson told BuzzFeed News that "digital materials are less expensive and a


5


good investment" that offer new features, like audio texts, personalized knowledge checks and expert videos. Its digital course materials save students up to 60\% compared to traditional printed textbooks, the company added. McGraw Hill didn't respond to a request for comment, but its CEO David Levin told the Financial Times in August that "in higher education, the era of the printed textbook is now over."


I) The textbook industry insists the online systems represent a better deal for students. "These digital products aren't just mechanisms for students to submit homework, they offer all kinds of features," David Anderson, the executive director of higher education with the Association of American Publishers, told BuzzFeed News. "It helps students understand in a way that you can't do with print homework assignments."


J) David Hunt, an associate professor in sociology at Augusta University, which has rolled out digital textbooks across its math and psychology departments, told BuzzFeed News that he understands the utility of using systems that require access codes. But he doesn't require his students to buy access to a learning program that controls the class assignments. "I try to make things as inexpensive as possible," said Hunt, who uses free digital textbooks for his classes but designs his own curriculum. "The online systems may make my life a lot easier but I feel like I'm giving up control. The discussions are the things where my expertise can benefit the students most."


K) A 20-year-old junior at Georgia Southern University told BuzzFeed News that she normally spends \$ 500-\$ 600 on access codes for class. In one case, the professor didn't require students to buy a textbook, just an access code to turn in homework. This year she said she spent \$ 900 on access codes to books and programs. "That's two months of rent," she said. "You can't sell any of it back. With a traditional textbook you can sell it for \$ 30 - \$ 50 and that helps to pay for your new semester's books. With an access code, you're out of that money. "


L) Benjamin Wolverton, a 19-year-old student at the University of South Carolina, told BuzzFeed News that "it's ridiculous that after paying tens of thousands in tuition we have to pay for all these access codes to do our homework." Many of the access codes he's purchased have been required simply to complete homework or quizzes. "Often it's only 10\% of your grade in class." he said. "You're paying so much money for something that hardly affects your grade—but if you didn't have it, it would affect your grades enough. It would be bad to start out at a B or C." Wolverton said he spent \$ 500 on access codes for digital books and programs this semester.


M) Harper, a poultry (家禽) science major, is taking chemistry again this year and had to buy a new access code to hand in her homework. She rented her economics and statistics textbooks for about \$ 20 each. But her access codes for homework, which can't be rented or bought second-hand, were her most expensive purchases: \$ 120 and \$ 85.


6


N) She still remembers the sting of her first experience skipping an assignment due to the high prices. "We don't really have a missed assignment policy," she said. "If you miss it, you just miss it. I just got zeros on a couple of first assignments. I managed to pull everything back up. But as a scared freshman looking at their grades, it's not fun."


36. A student's yearly expenses on access codes may amount to their rent for two months.


37. The online access codes may be seen as a way to tie the students to the digital system.


38. If a student takes a course again, they may have to buy a new access code to submit their


assignments.


39. McGraw Hill accounts for over one-fifth of the market share of college textbooks.


40. Many traditional textbook publishers are now offering online digital products, which they


believe will be the future of the publishing business.


41. One student complained that they now had to pay for access codes in addition to the high


tuition.


42. Digital materials can cost students less than half the price of traditional printed books


according to a publisher.


43. One student decided not to buy her access code until she received the pay for her part-time job.


44. Online systems may deprive teachers of opportunities to make the best use of their expertise


for their students.


45. Digital access codes are criticized because they are profit-driven just like the textbook business.


\subsection*{\textbf{Section C}}


\textbf{Directions:}\textit{ There are 2 passages in this section. Each passage is followed by some questions}\textit{ }\textit{or unfinished statements. For each of them there are four choices marked A), B), C) and}\textit{ D)}\textit{. You should decide on the best choice and mark the corresponding letter}\textbf{\textit{ on Answer Sheet2}} \textit{with a single line through the centre.}


\subsection*{\textbf{Passage One}}


\textbf{Questions 46 and 50 are based on the following passage.}


\subsection*{Passage One}


Questions 46 to 50 are based on the following passage.


Losing your ability to think and remember is pretty scary. We know the risk of dementia (痴呆症) increases with age. But if you have memory slips, you probably needn't worry. There are pretty clear differences between signs of dementia and age-related memory loss.


After age 50, it's quite common to have trouble remembering the names of people, places and things quickly, says Dr. Kirk Daffner of Brigham and Women's Hospital in Boston.


The brain ages just like the rest of the body. Certain parts shrink, especially areas in the brain that are important to learning, memory and planning. Changes in brain cells can affect communication between different regions of the brain. And blood flow can be reduced as blood vessels narrow.


7


Forgetting the name of an actor in a favorite movie, for example, is nothing to worry about. But if you forget the plot of the movie or don't remember even seeing it, that's far more concerning, Daffner says.


When you forget entire experiences, he says, that's "a red flag that something more serious may be involved." Forgetting how to operate a familiar object like a microwave oven, or forgetting how to drive to the house of a friend you've visited many times before can also be signs of something going wrong.


But even then, Daffner says, people shouldn't panic. There are many things that can cause confusion and memory loss, including health problems like temporary stoppage of breathing during sleep, high blood pressure, or depression, as well as medications ( 药物) like antidepressants.


You don't have to figure this out on your own. Daffner suggests going to your doctor to check on medications, health problems and other issues that could be affecting memory. And the best defense against memory loss is to try to prevent it by building up your brain's cognitive (认知的) reserve, Daffner says.


"Read books, go to movies, take on new hobbies or activities that force one to think in novel ways," he says. In other words, keep your brain busy and working. And also get physically active, because exercise is a known brain booster.


46.Why does the author say that one needn't be concerned about memory slips?


\begin{itemize}[leftmargin=2em,itemsep=.2em,topsep=.4em]\item[A.] Not all of them are symptoms of dementia.
\item[B.] They occur only among certain groups of people.
\item[C.] Not all of them are related to one's age.
\item[D.] They are quite common among fifty-year-olds.
\end{itemize}

47.What happens as we become aged according to the passage?


\begin{itemize}[leftmargin=2em,itemsep=.2em,topsep=.4em]\item[A.] Our interaction skills deteriorate.
\item[B.] Some parts of our brain stop functioning.
\item[C.] Communication within our brain weakens.
\item[D.] Our whole brain starts shrinking.
\end{itemize}

48.Which memory-related symptom should people take seriously?


\begin{itemize}[leftmargin=2em,itemsep=.2em,topsep=.4em]\item[A.] Totally forgetting how to do one's daily routines.
\item[B.] Inability to recall details of one's life experiences.
\item[C.] Failure to remember the names of movies or actors.
\item[D.] Occasionally confusing the addresses of one's friends.
\end{itemize}

49.What should people do when signs of serious memory loss show up?


\begin{itemize}[leftmargin=2em,itemsep=.2em,topsep=.4em]\item[A.] Check the brain's cognitive reserve.
\item[B.] Stop medications affecting memory.
\item[C.] Turn to a professional for assistance.
\item[D.] Exercise to improve their well-being.
\end{itemize}

50.What is Dr. Daffner's advice for combating memory loss?


8


\begin{itemize}[leftmargin=2em,itemsep=.2em,topsep=.4em]\item[A.] Having regular physical and mental checkups.
\item[B.] Taking medicine that helps boost one's brain.
\item[C.] Engaging in known memory repair activities.
\item[D.] Staying active both physically and mentally.Passage Two
\end{itemize}

\subsection*{\textbf{Passage Two}}


\textbf{Questions 51 to 55 are based on the following passage.}


A letter written by Charles Darwin in 1875 has been returned to the Smithsonian Institution Archives (档案馆) by the FBI after being stolen twice.


"We realized in the mid-1970s that it was missing," says Effie Kapsalis, head of the Smithsonian Insitution Archives. "It was noted as missing and likely taken by an intern (实习生), from what the FBI is telling us. Word got out that it was missing when someone asked to see the letter for research purposes," and the intern put the letter back. "The intern likely took the letter again once nobody was watching it."


Decades passed. Finally, the FBI received a tip that the stolen document was located very close to Washington, D.C. Their art crime team recovered the letter but were unable to press charges because the time of limitations had ended. The FBI worked closely with the Archives to determine that the letter was both authentic and definitely Smithsonian's property.


The letter was written by Darwin to thank an American geologist, Dr. Ferdinand Vandeveer Hayden, for sending him copies of his research into the geology of the region that would become Yellowstone National Park.


The letter is in fairly good condition, in spite of being out of the care of trained museum staff for so long. "It was luckily in good shape," says Kapsalis, "and we just have to do some minor things in order to be able to unfold it. It has some glue on it that has colored it slightly, but nothing that will prevent us from using it. After it is repaired, we will take digital photos of it and that will be available online. One of our goals is to get items of high research value or interest to the public online."


It would now be difficult for an intern, visitor or a thief to steal a document like this. "Archiving practices have changed greatly since the 1970s," says Kapsalis, "and we keep our high value documents in a safe that I don't even have access to."


51.What happened to Darwin's letter in the 1970s?


\begin{itemize}[leftmargin=2em,itemsep=.2em,topsep=.4em]\item[A.] It was recovered by the FBI.
\item[B.] It was stolen more than once.
\item[C.] It was put in the archives for research purposes.
\item[D.] It was purchased by the Smithsonian Archives.
\end{itemize}

52.What did the FBI do after the recovery of the letter?


\begin{itemize}[leftmargin=2em,itemsep=.2em,topsep=.4em]\item[A.] They proved its authenticity.
\item[B.] They kept it in a special safe.
\end{itemize}

9


\begin{itemize}[leftmargin=2em,itemsep=.2em,topsep=.4em]\item[C.] They arrested the suspect immediately.
\item[D.] They pressed criminal charges in vain.
\end{itemize}

53.What is Darwin's letter about?


\begin{itemize}[leftmargin=2em,itemsep=.2em,topsep=.4em]\item[A.] The evolution of Yellowstone National Park.
\item[B.] His cooperation with an American geologist.
\item[C.] Some geological evidence supporting his theory.
\item[D.] His acknowledgement of help from a professional.
\end{itemize}

54.What will the Smithsonian Institution Archives do with the letter according to Kapsalis?


\begin{itemize}[leftmargin=2em,itemsep=.2em,topsep=.4em]\item[A.] Reserve it for research purposes only.
\item[B.] Turn it into an object of high interest.
\item[C.] Keep it a permanent secret.
\item[D.] Make it available online.
\end{itemize}

55.What has the past half century witnessed according to Kapsalis?


\begin{itemize}[leftmargin=2em,itemsep=.2em,topsep=.4em]\item[A.] Growing interest in rare art objects.
\item[B.] Radical changes in archiving practices.
\item[C.] Recovery of various missing documents.
\item[D.] Increases in the value of museum exhibits.
\end{itemize}

\subsection*{\textbf{Part IV} \textbf{Translation} \textbf{(30 minutes)}}


\textbf{Directions:}\textit{ For this part, you are allowed 30minutes to translate a passage from Chinese into} \textit{English. You should write your answer}\textbf{\textit{ on Answer Sheet 2}}\textit{.}


过去，乘飞机出行对大多数中国人来说是难以想象的。如今，随着经济的发展和生活水平的提高，越来越多的中国人包括许多农民和外出务工人员都能乘飞机出行。他们可以乘飞机到达所有大城市，还有很多城市也在筹建机场。航空服务不断改进，而且经常会有廉价机票。 近年来，节假日期间选择乘飞机外出旅游的人数在不断增加。


10

\end{document}